\subsubsection{LightGBM}

Como hemos comentado antes, a diferencia de Random Forest, que construye árboles de forma independiente y paralela, LightGBM los construye secuencialmente, donde cada nuevo árbol se especializa en corregir los errores de los anteriores.

\subsubsection*{Justificación del modelo}

Hemos usado LightGBM debido a que presenta ventajas significativas para la tarea de detección de deepfakes de audio, como pueden ser su eficiencia computacional para grandes volúmenes de datos o su regulación itegrada con los parámetros L1 y L2.

\subsubsection*{Configuración experimental}

El proceso experimental para LightGBM se estructuró en cuatro fases: (1) establecimiento de un modelo baseline, (2) validación cruzada para evaluar robustez, (3) optimización de hiperparámetros mediante búsqueda sistemática, y (4) evaluación comparativa con métricas exhaustivas.

\paragraph{Preparación de datos}
A diferencia de otros algoritmos, LightGBM puede manejar directamente datos sin normalización, aunque requiere que no existan valores nulos.

\paragraph{Modelo baseline con LightGBM}
Se estableció una configuración inicial utilizando la interfaz nativa de LightGBM, que ofrece un control más fino sobre el proceso de entrenamiento:

\begin{lstlisting}[language=Python, caption=Configuración del modelo baseline LightGBM]
 # Creacion de datasets optimizados de LightGBM
 train_data = lgb.Dataset(X_train, label=y_train)
 test_data = lgb.Dataset(X_test, label=y_test, reference=train_data)

 # Parametros baseline
 params_baseline = {
     'objective': 'binary',              # Clasificacion binaria
     'metric': 'binary_logloss',          # Metrica de evaluacion
     'boosting_type': 'gbdt',              # Gradient Boosting Decision Tree
     'num_leaves': 31,                     # Numero maximo de hojas
     'learning_rate': 0.05,                 # Tasa de aprendizaje (eta)
     'feature_fraction': 0.9,               # Fraccion de caracteristicas por iteracion
     'bagging_fraction': 0.8,               # Fraccion de datos por iteracion
     'bagging_freq': 5,                      # Frecuencia de bagging
     'verbose': 0,                           # Modo silencioso
     'random_state': 12,
     'is_unbalance': True,                    # Compensacion de clases desbalanceadas
     'num_threads': -1                        # Usar todos los cores
 }

 # Entrenamiento con early stopping
 model_baseline = lgb.train(
     params_baseline,
     train_data,
     valid_sets=[test_data],
     num_boost_round=1000,
     callbacks=[
         lgb.early_stopping(50),   # Detiene si no mejora en 50 rondas
         lgb.log_evaluation(100)    # Log cada 100 iteraciones
     ]
 )
\end{lstlisting}

El parámetro \texttt{is\_unbalance=True} activa la compensación automática para clases desbalanceadas, ajustando los pesos de forma inversamente proporcional a la frecuencia de cada clase.

\paragraph{Early stopping}
Una ventaja importante de LightGBM es la capacidad de detener el entrenamiento cuando añadir más árboles no mejora el rendimiento en validación. Esto:
\begin{itemize}
    \item Previene el sobreajuste (overfitting)
    \item Ahorra tiempo computacional
    \item Determina automáticamente el número óptimo de iteraciones
\end{itemize}

\subsubsection*{Validación cruzada}

Para LightGBM se utilizó la interfaz compatible con scikit-learn (\texttt{LGBMClassifier}) que permite integrarse con las herramientas de validación cruzada estándar:

 \begin{lstlisting}[language=Python, caption=Validación cruzada con LightGBM]
 from sklearn.model_selection import StratifiedKFold, cross_val_score

 # Usamos la interfaz scikit-learn de LightGBM
 lgb_sklearn = lgb.LGBMClassifier(
     objective='binary',
     boosting_type='gbdt',
     num_leaves=31,
     learning_rate=0.05,
     random_state=12,
     is_unbalance=True,
     n_jobs=-1
 )

 cv = StratifiedKFold(n_splits=5, shuffle=True, random_state=12)

 cv_scores_accuracy = cross_val_score(lgb_sklearn, X_train, y_train, 
                                      cv=cv, scoring='accuracy')
 cv_scores_f1 = cross_val_score(lgb_sklearn, X_train, y_train, 
                                cv=cv, scoring='f1')
 cv_scores_auc = cross_val_score(lgb_sklearn, X_train, y_train, 
                                 cv=cv, scoring='roc_auc')
 \end{lstlisting}

La validación cruzada proporciona una estimación más robusta del rendimiento, reportándose la media y la desviación estándar de las métricas.

\subsubsection*{Optimización de hiperparámetros}

LightGBM cuenta con un conjunto de hiperparámetros específicos, diferentes a los de Random Forest. La búsqueda se centró en aquellos que más influyen en el rendimiento:

\begin{itemize}
    \item \texttt{num\_leaves}: Controla la complejidad del modelo. Valores altos permiten capturar patrones más complejos pero aumentan el riesgo de sobreajuste.
    
    \item \texttt{learning\_rate}: Tasa de aprendizaje (eta). Valores pequeños requieren más iteraciones pero pueden lograr mejor convergencia.
    
    \item \texttt{n\_estimators}: Número de iteraciones de boosting.
    
    \item \texttt{min\_child\_samples}: Mínimo de muestras requeridas en una hoja. Controla la profundidad y evita hojas con pocas muestras.
    
    \item \texttt{subsample} y \texttt{colsample\_bytree}: Fracciones de muestras y características para cada árbol, introduciendo aleatoriedad.
    
    \item \texttt{reg\_alpha} y \texttt{reg\_lambda}: Regularización L1 y L2 para controlar la complejidad.
\end{itemize}

\begin{lstlisting}[language=Python, caption=Grid Search para LightGBM]
 param_grid = {
     'num_leaves': [15, 31, 63],           # Complejidad del modelo
     'learning_rate': [0.01, 0.05, 0.1],    # Tasa de aprendizaje
     'n_estimators': [100, 200, 300],       # Numero de iteraciones
     'min_child_samples': [20, 50, 100],     # Minimo muestras por hoja
     'subsample': [0.6, 0.8, 1.0],           # Fraccion de muestras
     'colsample_bytree': [0.6, 0.8, 1.0],    # Fraccion de caracteristicas
     'reg_alpha': [0, 0.1, 0.5],             # Regularizacion L1
     'reg_lambda': [0, 0.1, 0.5]             # Regularizacion L2
 }

 grid_search = GridSearchCV(
     lgb.LGBMClassifier(objective='binary', random_state=42, 
                       is_unbalance=True, n_jobs=-1),
     param_grid,
     cv=cv,
     scoring='f1',
     n_jobs=-1,
     verbose=1
 )

 grid_search.fit(X_train, y_train)
 lgb_optimized = grid_search.best_estimator_
\end{lstlisting}

\subsubsection*{Métricas de evaluación}

Se emplearon las mismas métricas que en Random Forest para mantener la comparabilidad:

\begin{itemize}
    \item \textbf{Accuracy}: Proporción de predicciones correctas.
    \item \textbf{Precision}: Proporción de predicciones positivas correctas.
    \item \textbf{Recall}: Proporción de positivos reales correctamente identificados.
    \item \textbf{F1-Score}: Media armónica de precisión y recall.
    \item \textbf{AUC-ROC}: Área bajo la curva ROC.
    \item \textbf{MCC}: Coeficiente de correlación de Matthews.
\end{itemize}

La función de evaluación implementada es la misma que para Random Forest.